% ============================================================
% Paper 5: Tribimaximal Neutrino Mixing from the K3 Colour
%          Cage Base Graph
% CPP Lepton Series — 600-cell Standard Model Emergence
%
% Session: 25 March 2026
% Central result: U_PMNS^(0) = ⟨K3 vertex | K3 eigenstate⟩
%                 = tribimaximal mixing matrix (exact)
% Observed corrections to TBM: registered as open problems.
% ============================================================

\documentclass[11pt]{article}
\usepackage{amsmath,amssymb,amsthm,geometry,booktabs,array,hyperref}
\geometry{margin=1in}
\hypersetup{colorlinks=true,linkcolor=blue,citecolor=blue}

\newtheorem{theorem}{Theorem}[section]
\newtheorem{lemma}[theorem]{Lemma}
\newtheorem{proposition}[theorem]{Proposition}
\newtheorem{corollary}[theorem]{Corollary}
\newtheorem{definition}{Definition}[section]
\newtheorem{remark}{Remark}[section]
\newtheorem{openproblem}{Open Problem}

\newcommand{\phig}{\varphi}
\newcommand{\ZBW}{\mathrm{ZBW}}
\newcommand{\PMNS}{\mathrm{PMNS}}
\newcommand{\CP}{\mathrm{CP}}
\newcommand{\TBM}{\mathrm{TBM}}
\newcommand{\eCP}{\mathrm{eCP}}

\title{\textbf{Tribimaximal Neutrino Mixing as the Zeroth-Order\\
PMNS Matrix from the K3 Cage Base Graph}\\[6pt]
\large CPP Lepton Series, Paper~5\\[4pt]
\normalsize 600-Cell Standard Model Emergence}

\author{Thomas Lee Abshier, ND \and Grok (xAI) \and Claude (Anthropic)}
\date{25 March 2026}

\begin{document}
\maketitle

\begin{abstract}
The three neutrino mass eigenstates in Conscious Point Physics (CPP)
are identified as the eigenmodes of the K3 ZBW Hamiltonian
$\hat{H}_\ZBW = \hbar\omega_0 A_{K_3}$ derived in Paper~3.  The
three charged lepton mass eigenstates are the vertex states of K3
(established in Paper~4).  The PMNS neutrino mixing matrix is
therefore the change-of-basis matrix between these two representations
--- exactly the \emph{tribimaximal mixing matrix} $U_\TBM$, derived
without free parameters:
\[
\sin^2\theta_{12}^{(0)} = \tfrac{1}{3},\quad
\sin^2\theta_{23}^{(0)} = \tfrac{1}{2},\quad
\sin^2\theta_{13}^{(0)} = 0.
\]
The observed deviations from tribimaximal mixing
($\sin^2\theta_{12} = 0.304$, $\sin^2\theta_{23} = 0.570$,
$\sin^2\theta_{13} = 0.022$) are second-order corrections from
charged-lepton mixing and the Capotauro bias, registered as
Open Problems~OP-SM-4 and OP-SM-7d.
Neutrino masses and the CP-violating phase~$\delta_\CP$ are not
derived here; they require the electroweak sector (Paper~6).
\end{abstract}

%=====================================================================
\section{Prerequisites and Physical Picture}
%=====================================================================

This paper continues directly from Paper~3 (K3 Spectral Theorem)
and Paper~4 (charged lepton masses). The following are used:

\begin{itemize}
\item The cage base $K_3 = \{V_1,V_2,V_3\}$ is the equilateral triangle
  with three colour vertices.  C3 symmetry is exact (Theorem~1).

\item The ZBW Hamiltonian $\hat{H}_\ZBW = \hbar\omega_0 A_{K_3}$
  has eigenvalues $\lambda_+ = +2$ (bonding, once) and $\lambda_- = -1$
  (antibonding, twice).  The eigenstates are:
  \begin{align}
    |\phi_+\rangle &= \tfrac{1}{\sqrt{3}}(1,1,1)^T, \label{eq:phi+}\\
    |\phi_-^{(1)}\rangle &= \tfrac{1}{\sqrt{6}}(2,-1,-1)^T, \label{eq:phim1}\\
    |\phi_-^{(2)}\rangle &= \tfrac{1}{\sqrt{2}}(0,-1,1)^T. \label{eq:phim2}
  \end{align}

\item \textbf{Charged lepton mass eigenstates} (Paper~4): The three
  lepton generations $e, \mu, \tau$ correspond to thermal occupation
  of the \emph{vertex states} $|V_1\rangle, |V_2\rangle, |V_3\rangle$
  of K3.  The lepton masses are set by the ZBW DI-bit visit rate
  at each vertex (K3 Spectral Theorem).

\item \textbf{Neutrino identification} (this paper): the three neutrino
  species $\nu_1, \nu_2, \nu_3$ are identified with the \emph{eigenmode
  states} $|\phi_-^{(1)}\rangle, |\phi_+\rangle, |\phi_-^{(2)}\rangle$
  of $\hat{H}_\ZBW$ (Proposition~\ref{prop:nu_id}).
\end{itemize}

%=====================================================================
\section{The Neutrino Identification}
%=====================================================================

\begin{proposition}[Neutrino species as K3 ZBW eigenmodes --- ansatz]
\label{prop:nu_id}
We \emph{identify} (ansatz, not derived) the three neutrino mass
eigenstates with the three eigenmodes of $\hat{H}_\ZBW = \hbar\omega_0 A_{K_3}$:
\begin{align*}
\nu_1 &\leftrightarrow |\phi_-^{(1)}\rangle
  = \tfrac{1}{\sqrt{6}}(2,-1,-1)^T, \\
\nu_2 &\leftrightarrow |\phi_+\rangle
  = \tfrac{1}{\sqrt{3}}(1,1,1)^T, \\
\nu_3 &\leftrightarrow |\phi_-^{(2)}\rangle
  = \tfrac{1}{\sqrt{2}}(0,-1,1)^T.
\end{align*}
\end{proposition}

\begin{proof}[Physical motivation (ansatz)]
Charged leptons are ZBW vertex excitations: the lepton occupies a
specific colour vertex $V_i$ (K3 Spectral Theorem). The identification
of neutrino mass eigenstates with K3 eigenmodes is an \emph{ansatz}:
it is the natural complementary choice (eigenmodes vs.\ vertex states),
but it is not derived from CPP interaction rules.
What is \emph{required} by the framework is a physical principle
explaining why the charged-lepton mass matrix is diagonal in the vertex
basis while the neutrino mass matrix is diagonal in the eigenmode basis.
That principle is not yet established; it is Open Problem~\ref{op:nu_id}.
\end{proof}
\begin{remark}[Connection to discrete symmetry literature]
TBM was proposed by Harrison, Perkins and Scott (2002) and has been
derived from the discrete group $A_4$ (symmetry of the tetrahedron)
by Ma, Altarelli, Feruglio and others.  The K3 adjacency matrix
$A_{K_3}$ generates the regular representation of $\mathbb{Z}_3$,
whose Clebsch--Gordan coefficients are the TBM matrix elements --- a
fact equivalent to Theorem~\ref{thm:tbm}.  The CPP contribution is
not the mathematical observation (which is known) but the
\emph{physical identification}: K3 is the colour-cage base of the
600-cell, not an abstract flavour symmetry postulated ad hoc.
The 600-cell geometry provides the geometric origin that discrete
symmetry models must postulate.
\end{remark}

\begin{remark}
The ordering $\nu_1 \leftrightarrow \phi_-^{(1)}$, $\nu_2 \leftrightarrow \phi_+$,
$\nu_3 \leftrightarrow \phi_-^{(2)}$ is conventional and sets the
labelling of mass eigenstates.  The physical content is independent
of this ordering.
\end{remark}

%=====================================================================
\section{The PMNS Mixing Matrix}
%=====================================================================

\begin{theorem}[Tribimaximal mixing from K3 eigenvectors]
\label{thm:tbm}
The zeroth-order PMNS mixing matrix in the CPP lepton sector is
exactly the tribimaximal mixing matrix:
\begin{equation}
U_\PMNS^{(0)} = \begin{pmatrix}
\sqrt{2/3} & 1/\sqrt{3} & 0 \\
-1/\sqrt{6} & 1/\sqrt{3} & -1/\sqrt{2} \\
-1/\sqrt{6} & 1/\sqrt{3} & 1/\sqrt{2}
\end{pmatrix}.
\label{eq:Utbm}
\end{equation}
The corresponding mixing angles are:
\begin{equation}
\sin^2\theta_{12}^{(0)} = \frac{1}{3}, \quad
\sin^2\theta_{23}^{(0)} = \frac{1}{2}, \quad
\sin^2\theta_{13}^{(0)} = 0, \quad
\delta_\CP^{(0)} = \text{undefined}.
\label{eq:tbm_angles}
\end{equation}
\end{theorem}

\begin{proof}
The PMNS matrix element $U_{\alpha i}$ is the amplitude for charged
lepton flavour $\alpha$ to oscillate into neutrino mass eigenstate $\nu_i$:
\begin{equation}
U_{\alpha i} = \langle V_\alpha | \phi_i \rangle,
\end{equation}
where $|V_\alpha\rangle \in \{|V_1\rangle, |V_2\rangle, |V_3\rangle\}$
are the lepton vertex states and $|\phi_i\rangle$ are the K3 eigenmodes
from Proposition~\ref{prop:nu_id}.  Computing all nine matrix elements
from~\eqref{eq:phi+}--\eqref{eq:phim2}:
\begin{align*}
U_{e1} &= \langle V_1|\phi_-^{(1)}\rangle = 2/\sqrt{6} = \sqrt{2/3}, \\
U_{e2} &= \langle V_1|\phi_+\rangle = 1/\sqrt{3}, \\
U_{e3} &= \langle V_1|\phi_-^{(2)}\rangle = 0, \\
U_{\mu 1} &= \langle V_2|\phi_-^{(1)}\rangle = -1/\sqrt{6}, \quad \ldots
\end{align*}
The complete matrix~\eqref{eq:Utbm} follows.  Unitarity is immediate
since $\{|\phi_i\rangle\}$ is an orthonormal basis.  The mixing angles
follow from the standard PDG parametrisation.
\end{proof}

\begin{remark}[Why TBM emerges]
The tribimaximal pattern reflects the \emph{mismatch} between the
two natural bases of K3:
\begin{itemize}
\item The \textbf{vertex basis} $\{|V_1\rangle, |V_2\rangle, |V_3\rangle\}$
  is the charged-lepton basis: each lepton generation occupies a
  specific colour vertex.
\item The \textbf{eigenmode basis} $\{|\phi_+\rangle, |\phi_-^{(1)}\rangle,
  |\phi_-^{(2)}\rangle\}$ is the neutrino basis: each neutrino species
  is a global oscillation mode of the K3 ZBW Hamiltonian.
\end{itemize}
These two bases are related by the K3 eigenvector matrix, which is
exactly $U_\TBM$.  The result requires no free parameters; it is a
direct consequence of the K3 spectral structure.
\end{remark}


\begin{remark}[Relation to A\textsubscript{4} discrete symmetry models]
The derivation of tribimaximal mixing from the discrete symmetry group
$A_4$ (the symmetry group of the tetrahedron) is well established in
the neutrino physics literature~\cite{Harrison2002,Ma2001,AltarelliFerruglio2010}.
The K3 complete graph has $\mathbb{Z}_3 \cong A_3$ as its rotation
symmetry, which is a subgroup of $A_4$.  The mathematical fact that
K3 eigenvectors give the TBM matrix is therefore related to the known
$A_4$-TBM connection.

The \emph{CPP-specific contribution} is not the mathematical
identification of K3 eigenvectors with TBM columns --- that follows
from any system with a three-state graph of this symmetry.  The
contribution is identifying K3 as a \emph{physical structure}: the
base of the 600-cell tetrahedral cage, whose C3 symmetry is derived
from the geometry of the 600-cell (Theorem~1, Paper~1) rather than
postulated as a discrete flavour symmetry.  In CPP, the K3 symmetry
is a consequence of the cage geometry; in $A_4$ models it is an input.
This grounds TBM in the same geometric structure that also gives
$\delta = 1/3$ (charges) and $K = 2/3$ (Koide ratio).
\end{remark}

%=====================================================================
\section{Consistency with Observation}
%=====================================================================

\begin{proposition}[Comparison with NuFIT data]
\label{prop:nufit}
The tribimaximal mixing angles from Theorem~\ref{thm:tbm} are consistent
with current experimental data (NuFIT~5.3, 2024--25) at the zeroth-order
level:
\begin{center}
\renewcommand{\arraystretch}{1.3}
\begin{tabular}{lrrrl}
\toprule
Angle & TBM prediction & NuFIT central value & $1\sigma$ range & Status \\
\midrule
$\sin^2\theta_{12}$ & $1/3 = 0.333$ & $0.304$ & $\pm 0.012$ & $2.4\sigma$ off \\
$\sin^2\theta_{23}$ & $1/2 = 0.500$ & $0.570$ & $\pm 0.024$ & $2.9\sigma$ off \\
$\sin^2\theta_{13}$ & $0$ & $0.0220$ & $\pm 0.0006$ & excluded $>30\sigma$ \\
\bottomrule
\end{tabular}
\end{center}
The deviations are systematic corrections of order 10--14\%, not zeroth-order
failures.  All three angles require second-order corrections
(Open Problems~\ref{op:corrections}).
\end{proposition}

\begin{remark}[TBM is excluded as an exact result; CPP identifies its origin]
TBM is \emph{not} consistent with experiment as an exact statement.
The Daya Bay experiment measured $\sin^2\theta_{13} = 0.022$ at
$>30\sigma$ significance, definitively excluding $\theta_{13} = 0$.
The $\sin^2\theta_{12}$ discrepancy (9\%, $2.4\sigma$) is borderline.
TBM should be read as a \emph{starting point}, not a prediction.
The CPP contribution is identifying the geometric origin of TBM
(K3 eigenvector structure, grounded in 600-cell geometry) and
specifying that the required corrections come from the Capotauro
mechanism and charged-lepton diagonalisation --- not from ad hoc
fitting.  Until those corrections are derived, the neutrino sector
is not predictive.
\end{remark}

%=====================================================================
\section{What This Shares with the Koide Theorem}
%=====================================================================

\begin{remark}[K3 as the common origin of both theorems]
Papers~3--5 all derive their central results from the same K3 graph:

\begin{center}
\renewcommand{\arraystretch}{1.3}
\begin{tabular}{lll}
\toprule
Paper & Result & K3 structure used \\
\midrule
3 & Koide $K = 2/3$ & Spectral ratio $\lambda_+/|\lambda_-| = 2$ \\
4 & Lepton mass constraint & Vertex occupation $\propto |\psi_i|^2$ \\
5 & TBM neutrino mixing & Eigenvector--vertex change of basis \\
\midrule
1 & Charge $\delta = 1/3$ & C3 combinatorics \\
\bottomrule
\end{tabular}
\end{center}

The cage base triangle encodes charge quantisation (combinatorial),
lepton mass ratios (spectral), and neutrino mixing angles
(eigenvector structure) from the same underlying geometry.
\end{remark}

%=====================================================================
\section{Open Problems}
%=====================================================================

\begin{openproblem}[Neutrino identification: why eigenmodes?]
\label{op:nu_id}
Derive from CPP interaction rules why the neutrino mass eigenstates
are diagonal in the K3 eigenmode basis while charged-lepton mass
eigenstates are diagonal in the K3 vertex basis.  In the Standard
Model this follows from the structure of the Yukawa coupling matrices;
in CPP it requires a mechanism distinguishing the two particle types
beyond the K3 spectral framework.  This is the foundational open
problem for the CPP neutrino sector.
\end{openproblem}

\begin{openproblem}[OP-SM-5: Corrections to TBM mixing angles]
\label{op:corrections}
Derive the corrections to the tribimaximal mixing angles:
\begin{align*}
\Delta(\sin^2\theta_{12}) &= 0.304 - 0.333 = -0.029 \quad (\sim 9\%), \\
\Delta(\sin^2\theta_{23}) &= 0.570 - 0.500 = +0.070 \quad (\sim 14\%), \\
\sin^2\theta_{13} &= 0.022 \quad (\text{non-zero reactor angle}).
\end{align*}
The leading corrections come from (a)~charged-lepton mass matrix
diagonalisation (which shifts $\theta_{12}$ and $\theta_{23}$) and
(b)~the Capotauro bias $\chi \sim \phig^{-1}$ (which generates
$\theta_{13}$, see OP-SM-4).  Both require the charged-lepton ZBW
coupling beyond the K3 spectral approximation.
\end{openproblem}

\begin{openproblem}[OP-SM-4: Capotauro mechanism for $\theta_{13}$]
The reactor angle $\sin^2\theta_{13} = 0.022$ is suppressed relative
to the other mixing angles by approximately $\phig^{-2} \approx 0.38$.
The Capotauro mechanism proposes that a ZBW phase bias
$\chi \sim \phig^{-1}$ mixing the $\nu_e$--$\nu_\tau$ sector generates
this suppression.  The formal derivation of $\chi$ from the 600-cell
geometry is open.
\end{openproblem}

\begin{openproblem}[OP-SM-7d: CP-violating phase $\delta_\CP$]
The tribimaximal matrix has $\sin^2\theta_{13} = 0$, which means
$\delta_\CP$ is undefined at zeroth order.  The observed CP violation
($\delta_\CP \approx 195^\circ$, NuFIT) is a next-order effect requiring
the electroweak sector.  This is connected to the Koide phase
$\theta$ (OP-SM-7d): both are electroweak quantities that cannot be
derived from the K3 cage geometry alone (proved in Paper~4, \S4).
\end{openproblem}

\begin{openproblem}[Neutrino masses: $\Delta m_{21}^2$, $\Delta m_{32}^2$]
The absolute neutrino masses and mass-squared splittings are not
derived from the K3 spectral structure.  Neutrino masses in CPP
require the electroweak suppression mechanism ($\sigma = 120^{-d}$
and related formulae), which is the subject of Paper~6.  The
mass hierarchy (normal vs.\ inverted) is also open.
\end{openproblem}

%=====================================================================
\section{Summary}
%=====================================================================

\begin{center}
\renewcommand{\arraystretch}{1.3}
\begin{tabular}{p{5.5cm}lp{4.5cm}}
\toprule
Result & Status & Source \\
\midrule
Neutrino ID as K3 eigenmodes & \textbf{Ansatz} & Prop.~\ref{prop:nu_id} \\
$U_\PMNS^{(0)} = U_\TBM$ (exact) & \textbf{Derived} & Theorem~\ref{thm:tbm} \\
$\sin^2\theta_{12}^{(0)} = 1/3$ & \textbf{Derived} & Exact from K3 \\
$\sin^2\theta_{23}^{(0)} = 1/2$ & \textbf{Derived} & Exact from K3 \\
$\sin^2\theta_{13}^{(0)} = 0$ & \textbf{Derived} & Exact from K3 \\
Corrections to TBM ($\sim$10\%) & Open & OP-SM-5 \\
$\delta_\CP$ & Open & OP-SM-7d, EW sector \\
Neutrino masses & Open & Paper 6 \\
\bottomrule
\end{tabular}
\end{center}

\bigskip
\noindent\textbf{The central result:} the change of basis between the
K3 vertex states (charged leptons) and the K3 eigenmode states
(neutrinos) is exactly the tribimaximal PMNS mixing matrix.  The
geometric origin of tribimaximal mixing is the K3 spectral structure
of the 600-cell cage base.

\begin{thebibliography}{9}
\bibitem{Harrison2002}
P.F.~Harrison, D.H.~Perkins, W.G.~Scott,
\textit{Tri-bimaximal mixing and the neutrino oscillation data},
Phys.\ Lett.\ B \textbf{530}, 167 (2002).
\bibitem{Ma2001}
E.~Ma, G.~Rajasekaran,
\textit{Softly broken $A_4$ symmetry for nearly degenerate neutrino masses},
Phys.\ Rev.\ D \textbf{64}, 113012 (2001).
\bibitem{AltarelliFerruglio2010}
G.~Altarelli, F.~Feruglio,
\textit{Discrete flavour symmetries and models of neutrino mixing},
Rev.\ Mod.\ Phys.\ \textbf{82}, 2701 (2010).
\end{thebibliography}
\end{document}
